% ============================================
% 1 问题重述
% ============================================
\section{问题重述}

\subsection{问题背景}

古代玻璃是丝绸之路贸易往来的重要物证。我国古代玻璃在吸收西亚、埃及技术后，就地取材制作，外观与外来玻璃相似但化学成分不同。玻璃的主要原料为石英砂（$\text{SiO}_2$），需添加助熔剂降低熔点——不同助熔剂决定玻璃类型：铅钡玻璃以铅矿石为助熔剂（PbO、BaO 含量高），是我国自主发明的玻璃品种；钾玻璃以草木灰为助熔剂（$\text{K}_2\text{O}$ 含量高），流行于岭南及东南亚。

玻璃在埋藏环境中极易风化，内部元素与环境元素发生交换，导致成分比例变化，影响类别判断。

\subsection{数据说明}

\begin{table}[H]
    \centering
    \caption{数据文件说明}
    \label{tab:data}
    \begin{tabular}{p{5cm} p{5cm} p{4cm}}
        \toprule
        \textbf{数据文件} & \textbf{内容} & \textbf{规模} \\
        \midrule
        \texttt{表单1\_文物信息.csv} & 文物元数据（编号、纹饰、类型、颜色、风化状态） & 58 条 \\
        \texttt{表单2\_化学成分.csv} & 14 种氧化物的重量百分比 & 69 个采样点 \\
        \texttt{表单3\_未知类别.csv} & 待分类样品的化学成分 & 8 个样品（A1--A8） \\
        \bottomrule
    \end{tabular}
\end{table}

有效数据标准：各成分比例累加和在 $85\%\sim105\%$ 之间。

\subsection{需解决的问题}

\begin{enumerate}[label=\arabic*.]
    \item \textbf{风化分析}：分析表面风化与类型、纹饰、颜色的关系；结合类型分析风化前后化学成分统计规律；预测风化点风化前的化学成分含量。
    \item \textbf{分类与亚类划分}：分析高钾与铅钡玻璃的分类规律；选择合适的化学成分进行亚类划分，评价合理性与敏感性。
    \item \textbf{未知类别鉴别}：对 8 个未知样品进行类型鉴别，分析分类敏感性。
    \item \textbf{关联关系分析}：分析不同类别玻璃化学成分间的关联关系，比较类型间差异。
\end{enumerate}
